\chapter{Il Mercato Elettrico Italiano}
\label{cap:mercato}

Questo capitolo descrive l'architettura del mercato elettrico italiano limitatamente a
quanto serve per il lavoro: la struttura delle sedi di negoziazione, il funzionamento
dell'asta del Mercato del Giorno Prima e la natura della variabilit\`a dei prezzi. La
descrizione del MGP \`e basata sui dati effettivamente utilizzati, le cui caratteristiche
sono documentate nel Capitolo~\ref{cap:simulazione}.

\section{Struttura del GME (Gestore dei Mercati Energetici)}
\label{sec:gme}

Il Gestore dei Mercati Energetici (GME) organizza le sedi di negoziazione dell'energia
elettrica all'ingrosso in Italia. Ai fini di questo lavoro rilevano tre caratteristiche
della sua architettura.

La prima \`e la \textbf{sequenzialit\`a} delle sedi: il Mercato del Giorno Prima si chiude
prima dei Mercati Infragiornalieri, che a loro volta precedono il Mercato per il Servizio di
Dispacciamento. Un operatore che modifica la propria posizione lo fa quindi in momenti
successivi e con informazioni via via pi\`u precise.

La seconda \`e la \textbf{zonalit\`a}: il territorio nazionale \`e suddiviso in zone di
mercato, ciascuna con un proprio prezzo quando i limiti di transito fra zone sono
vincolanti. Nel giorno pilota analizzato compaiono sette zone fisiche italiane — NORD, CNOR,
CSUD, SUD, CALA, SICI, SARD — accanto a zone \emph{virtuali} che rappresentano le frontiere
con l'estero.

La terza \`e la \textbf{trasparenza dei dati}: il GME pubblica, per ciascun giorno e per
ciascuna sede, l'insieme completo delle offerte presentate, comprese quelle rifiutate. \`E
questa disponibilit\`a che rende possibile la ricostruzione delle curve aggregate su cui si
fonda l'intero lavoro: senza il dettaglio delle offerte non accettate la curva sarebbe
troncata proprio nel punto che interessa, cio\`e attorno all'equilibrio.

\section{MGP e MI (Mercato del Giorno Prima e Mercati Infragiornalieri)}
\label{sec:mgp}

\subsection*{Il meccanismo d'asta}

Il Mercato del Giorno Prima \`e un'\textbf{asta a prezzo uniforme}: gli operatori presentano
coppie prezzo--quantit\`a, distinte fra acquisto e vendita, riferite a un singolo periodo e a
una singola zona. L'algoritmo di mercato ordina le offerte di vendita per prezzo crescente e
quelle di acquisto per prezzo decrescente, ottenendo due funzioni a gradini, e ne individua
l'intersezione. Tutte le offerte accettate sono remunerate al medesimo prezzo, quello di
equilibrio, indipendentemente dal prezzo a cui erano state presentate.

Formalmente, indicando con $S(p)$ la quantit\`a complessivamente offerta in vendita a un
prezzo non superiore a $p$ e con $D(p)$ quella domandata a un prezzo non inferiore a $p$, il
prezzo di equilibrio \`e
\begin{equation}
p^{*} \;=\; \min \{\, p \;:\; S(p) \ge D(p) \,\}.
\label{eq:clearing}
\end{equation}

La funzione $S$ \`e non decrescente e la funzione $D$ non crescente, quindi l'eccesso di
offerta $S(p) - D(p)$ \`e monotono non decrescente e l'insieme dei prezzi che pareggiano il
mercato \`e sempre una semiretta: il prezzo di equilibrio ne \`e l'estremo inferiore. Quando
l'eccesso si annulla su un intervallo di prezzi l'equilibrio non \`e unico e la scelta del
punto diventa una convenzione; sui dati esaminati questo accade nello 0{,}8\% delle aste
orarie e in nessuna di quelle a quarto d'ora, quindi non costituisce una fonte rilevante di
ambiguit\`a.

\subsection*{Il periodo di mercato}

Il periodo elementare d'asta (\emph{Market Time Unit}) \`e cambiato durante il periodo
coperto dai dati. Fino al 30 settembre 2025 il mercato era organizzato su base
\textbf{oraria}, con ventiquattro aste indipendenti al giorno; dal 1\textsuperscript{o}
ottobre 2025 \`e passato al \textbf{quarto d'ora}, con novantasei aste al giorno. Il
passaggio \`e netto e verificabile nei dati: il 30 settembre 2025 la totalit\`a delle offerte
\`e oraria, il giorno successivo l'82{,}8\% \`e a quarto d'ora.

Questa discontinuit\`a ha conseguenze dirette sul disegno dello studio, discusse nella
Sezione~\ref{sec:dataset}: l'anno solare 2025 non \`e omogeneo, e la risoluzione temporale
non \`e neutrale per la valutazione di una strategia di arbitraggio, che dipende da quanto
\`e fine la griglia su cui l'accumulo pu\`o caricare e scaricare.

\subsection*{Tipologie di offerta}

Accanto alle offerte \emph{semplici}, riferite a un singolo periodo e frazionabili, esistono
le \textbf{offerte a blocchi}, accettabili solo per intero e su tutti i periodi che coprono.
Sui dati esaminati un blocco copre da quattro a novantasei periodi, ha un prezzo unico e uno
stato identico in tutti i suoi periodi: non si osserva alcun blocco accettato soltanto su una
parte di essi.

La presenza di offerte a blocchi rende il problema d'asta un programma con vincoli di
interezza e genera il fenomeno del \emph{rifiuto paradossale}: un blocco pu\`o risultare in
merito sul prezzo e venire comunque respinto, perch\'e accettarlo renderebbe inconsistente la
soluzione. I dati riportano questi casi con un codice di stato dedicato.

\subsection*{I Mercati Infragiornalieri}

% DA COMPLETARE: descrizione dei MI, da sviluppare se l'analisi verra' estesa oltre il MGP.
% Elementi da trattare: sessioni ad asta MI-A1/A2/A3, negoziazione continua XBID, funzione
% di aggiustamento delle posizioni assunte sul MGP, e implicazioni per un accumulo che
% volesse operare anche su queste sedi.

\section{MSD (Mercato per il Servizio di Dispacciamento)}
\label{sec:msd}

% DA COMPLETARE: descrizione del MSD.
% Elementi da trattare:
%   - finalita': approvvigionamento delle risorse per il bilanciamento in tempo reale
%     e per la risoluzione delle congestioni, con Terna come controparte unica;
%   - meccanismo pay-as-bid, diverso dal prezzo uniforme del MGP;
%   - rilevanza per un BESS: i ricavi da MSD sono spesso superiori a quelli da arbitraggio
%     sul MGP, quindi ignorarli sottostima la redditivita' complessiva dell'investimento.
%     Questa esclusione va dichiarata fra i limiti del modello (Sezione 6.2).

\section{L'aleatorietà del prezzo}
\label{sec:aleatorieta}

Il prezzo zonale \`e la variabile aleatoria centrale del problema. Due sue caratteristiche
sono direttamente rilevanti per la strategia di accumulo.

\paragraph{Ampiezza dell'escursione infragiornaliera.} \`E il differenziale fra ore a prezzo
basso e ore di picco, cio\`e esattamente il margine lordo di un ciclo di arbitraggio. Nel
mese di gennaio 2025 il prezzo zonale della zona NORD si \`e mosso fra un minimo orario medio
di circa 116~\euromwh{} nelle ore notturne e un massimo di circa 172~\euromwh{} nelle ore di
picco serale.

\paragraph{Limiti di prezzo e prezzi negativi.} I prezzi ammessi sul mercato non sono
confinati ai valori positivi: sui dati esaminati le offerte spaziano da $-500$ a
$4\,000$~\euromwh{}. La possibilit\`a di prezzi negativi \`e caratteristica dei sistemi ad
alta penetrazione rinnovabile, dove in alcune ore conviene pagare per immettere energia
piuttosto che modulare la produzione, ed \`e proprio la condizione in cui l'accumulo trova le
occasioni di acquisto pi\`u convenienti.

\paragraph{Rigidit\`a della domanda.} Una quota rilevante della domanda \`e presentata al
prezzo massimo, cio\`e \`e disposta ad acquistare a qualunque prezzo: nel giorno pilota si
tratta del 52\% delle offerte di acquisto della zona NORD. La curva di domanda risulta quindi
molto ripida e l'effetto di un accumulo sul prezzo passa quasi interamente attraverso la
curva di offerta. \`E un elemento da tenere presente nell'interpretazione degli scenari del
Capitolo~\ref{cap:simulazione}.

% DA COMPLETARE: analisi statistica della serie dei prezzi (stagionalita' settimanale e
% annuale, persistenza, distribuzione dell'escursione giornaliera). Da sviluppare quando
% sara' disponibile la ricostruzione su tutto l'anno 2025.
